\documentclass[10pt,twocolumn,letterpaper]{article}
\usepackage[margin=0.75in]{geometry}
\usepackage{amsmath,amssymb,amsthm,booktabs,hyperref}
\hypersetup{colorlinks=true,linkcolor=blue,citecolor=blue,urlcolor=blue}
\setlength{\columnsep}{0.27in}
\setlength{\emergencystretch}{2em}
\newtheorem{theorem}{Theorem}
\newtheorem{lemma}{Lemma}
\newtheorem{proposition}{Proposition}
\newcommand{\R}{\mathbb{R}}
\newcommand{\Q}{\mathbb{Q}}
\newcommand{\calD}{\mathcal{D}}
\newcommand{\calS}{\mathcal{S}}
\newcommand{\Qfix}{Q_{\mathrm{fix}}}
\newcommand{\pending}[1]{\par\noindent\textbf{Pending evidence.} #1\par}
\title{Certified Calibration Repair from Compressed Fixed-Anchor Features}
\author{Research working draft\\Author and affiliation details pending}
\date{10 October 2026\\\textbf{Working draft---not submission-ready}}
\begin{document}
\maketitle

\begin{abstract}
Removing a calibration record from a quantized language model raises two distinct
questions: which discrete model should replace it, and which persistent state is
needed for subsequent requests? We study a precisely specified counterfactual in
which a fixed pretrained checkpoint and a membership-independent nearest-grid
anchor determine source-local calibration features. The target interprets finite
feature words as exact dyadic numbers, retains the original normalization, and
uses fixed grids and deterministic adaptive rounding. We describe repair from
compressed feature enclosures, with universal decision certificates, bounded
exact replay, stage-specific certified solvers, and canonical successor states.
We prove conditional code and state correctness; acceptance, termination under a
resource cap, and practical benefit are separate obligations. Historical
single-stage development experiments show a tradeoff rather than universal
superiority: at 768 retained tokens, forty-bit compression saves 20.24\% of the
declared payload while taking 10.29\% longer than a lossless cache; at 1,536 tokens,
native exact Gram deletion is both faster and smaller. A complete-model pilot
matches all 42.47 million codes and canonical successor bytes over two deletions.
Compression takes 1.79 and 1.60 times the lossless-cache request time, while saving
less than 1\% of declared deployment storage including the checkpoint. This is a
bounded correctness result with an adverse cost tradeoff, not competitive
acceleration. On eight previously exposed articles, final-model perplexity is
53.72 versus 52.82 for matched sequential calibration: the 1.71\% regression
passes the frozen development tolerance but establishes no quality advantage.
Useful-scale validation, untouched quality evaluation, a second model,
prospective confirmation, and broader comparisons remain necessary before
submission.
\end{abstract}

\section{Introduction}
Post-training quantization uses calibration data to select low-bit model weights.
When calibration membership changes, an application may need a replacement that
matches an explicitly defined retained-data reconstruction. Repeating the entire
calibration computation is one option. Keeping all exact feature values is
another. A smaller archive is useful only if it preserves enough information to
recover the required discrete decisions, or can detect when additional source
access is needed.

Our question is therefore about the joint cost of \emph{preparing, retaining,
accessing, and updating calibration evidence}. A repair method is not useful
merely because it is faster than an unoptimized cold constructor. An equally
informed indexed constructor can execute the same repair procedure. Strong
comparisons must include exact cached features, exact additive statistics, and
fresh reconstruction under matching output and access contracts.

We examine compressed, source-local feature enclosures. Compression discards
exact feature words from the live descriptor while retaining intervals that
contain them. A numerical proposal suggests coefficients or codes; a separate
outward verifier must certify each decision for every realization of the feature
box. An unresolved stage can trigger an admitted exact replay and point
certificate. If those checks do not complete, the service refuses rather than
returning an approximate replacement.

The counterfactual is deliberately restricted. A fixed nearest-grid anchor,
constructed from the pretrained weights alone, produces every calibration
feature. Repaired weights do not change later calibration features. This differs
from ordinary sequential calibration and from stock floating-point GPTQ
implementations~\cite{gptq}. It also concerns calibration membership, not removal
of information from the checkpoint's pretraining.

The prospective contribution is exact discrete-code recovery from uncertain
archived evidence with useful complete costs. Fixed features, sufficient
statistics, matrix identities, and decision stability are classical ingredients;
we do not claim them as new. This draft separates the conditional correctness
argument from the still-incomplete evidence for practical and scientific value.
The target boundaries follow the V41 project handoff~\cite{boundaries};
the current draft adds the complete-model V43 and diagnostic V44 evidence~\cite{v43,v44}.

\section{An explicit reconstruction target}
\label{sec:target}
Let $W$ be a fixed pretrained checkpoint and let $\calD_0$ be an original set of
identified token records. Fix the tokenizer, record boundaries, stage order,
coordinate order, ridge parameters, row grids, and finite feature evaluator.
The original normalization $N_0>0$ is part of this specification; deletion does
not replace it with the number of surviving tokens.

For each stage $\ell$, construct a nearest-grid anchor matrix from the stage's
base weights, using lower-code midpoint ties. The full anchor depends only on
$W$ and these frozen grids. For source $j$, let
\[
 X_{\ell j}=F_\ell(W_{\mathrm{anchor}},j)
       \in\Q^{d_\ell\times t_j}.
\]
Here $F_\ell$ is the declared finite CPU computation, followed by exact
interpretation of each successful binary64 output word as a dyadic rational.
It is not an ideal real-arithmetic transformer. Unsupported or unresolved finite
operations cause refusal, so this evaluator is partial. Each source is evaluated
independently; its features cannot depend on another source's membership.

For a retained set $R\subseteq\calD_0$, concatenate its features in canonical
source order to obtain $X_\ell(R)$, and define
\begin{equation}
 H_\ell(R)=\lambda_\ell I+
   \frac{X_\ell(R)X_\ell(R)^\top}{N_0},
 \qquad \lambda_\ell>0.
 \label{eq:metric}
\end{equation}
A target with stage-specific original normalizations can be defined similarly;
all empirical configurations discussed here use one frozen token count.

\paragraph{Exact adaptive rounding.}
For one stage, write its rational positive-definite metric as a reverse LDL
factorization
\[
 H=L^\top\operatorname{diag}(\tau)L,
 \qquad B=I-L,
\]
where $L$ is unit lower triangular and every pivot $\tau_i$ is positive. For each
base-weight row $w$, process coordinates in their frozen order:
\begin{align}
 v_i&=w_i-\sum_{h<i} B_{ih}(w_h-q_h),\label{eq:recurrence}\\
 q_i&=\operatorname{nearest}_{G_i}(v_i).
 \label{eq:rounding}
\end{align}
Stage and output-row indices are suppressed in this notation. In V43, each
output row has one base-only dyadic24 scale $s_{\ell,r}$ and grid
$G_{\ell,r}=s_{\ell,r}\{-8,\ldots,7\}$, shared across its coordinates.
Each finite increasing rational grid $G_i$ is fixed before deletion. Nearest
rounding saturates at endpoints; an exact midpoint selects the lower code. The
collection of all stage outputs defines $\Qfix(R)$ wherever the required feature
evaluations succeed. Positive ridge makes the mathematical suffix systems
invertible, but does not guarantee that a bounded numerical implementation will
certify every decision.

\paragraph{What remains fixed.}
The checkpoint, anchor, original normalization, grids, arithmetic target, and
ordering belong to target identity. Changing any of them requires a new target.
Fresh reconstruction used as an oracle must retain the same original
normalization. Matching a separately retuned or sequentially recalibrated model
would answer a different question.

\section{Repair from compressed evidence}
\label{sec:method}
Trusted preparation evaluates each original source once through all stages.
For each $X_{\ell j}$ it constructs a canonical descriptor $D_{\ell j}$ whose
public decoding is an entrywise box
\begin{equation}
 \underline X_{\ell j}\leq X_{\ell j}\leq\overline X_{\ell j}.
 \label{eq:containment}
\end{equation}
Preparation checks this containment against the actual finite values. Descriptor
metadata binds the source ID, token identity, stage, target, codec settings, and
source-word commitment. A hash detects changed bytes; it does not independently
prove containment or authenticate an untrusted execution host.

A request validates the predecessor, retained membership, deletion IDs, and all
target and provenance bindings. Retained source-local descriptors remain
unchanged. At each stage the service presents their decoded boxes to a universal
certificate. The numerical solver may use floating-point proposals, but proposals
alone never establish correctness. Directed residual bounds and decision-cell
checks supply the evidence.

\paragraph{Decision cells.}
For $G_i=\{g_0<\cdots<g_{k-1}\}$, let
$m_a=(g_a+g_{a+1})/2$. The cell selecting an interior code $g_a$ is
$(m_{a-1},m_a]$; endpoint cells extend to infinity. Thus a closed enclosure
$[\underline v_i,\overline v_i]$ certifies $g_a$ only when it lies wholly within
that cell. A lower boundary is strict and an upper boundary is inclusive. This
asymmetry is required by the lower-code tie rule, including exact midpoint cases.
Previously certified codes are used when bounding the next adaptive decision.

\paragraph{Bounded fallback.}
If a box certificate refuses, the service may regenerate the retained sources'
exact finite features and invoke an admitted point certificate. Every traversed
ancestor is charged, including stages whose compressed certificates already
accepted. Replayed source words and their containment must match trusted
preparation. Repeated stage failures may reuse an already regenerated source
within the same request, with that access made explicit. Temporary exact features
do not silently replace the canonical compressed representation. Exhausted CPU,
memory, refinement, or replay allowances cause refusal without a complete model
release.

\subsection{Stage-specific numerical representations}
For width $d$, retained token count $T$, and output-row count $r$, token-space
coefficient construction has a leading structural term proportional to $dT^2$.
Feature-space construction has terms proportional to $d^2T+d^3$, followed by
row work proportional to $rd^2$. These expressions describe arithmetic schedules,
not measured runtime. Interval width, repeated verification, decoding, source
replay, and native implementation constants can reverse finite-size rankings.

The feature-space route builds directed Gram enclosures from actual feature
boxes, or accepts an exact Gram with trusted source lineage. It must not infer
positive semidefiniteness for every arbitrary matrix in an entrywise Gram box.
Every \emph{realized feature Gram} is positive semidefinite; that is the premise
used by the certificate.

For $K=XX^\top$, $\beta=\lambda N_0$, and a suffix $S$, let
$A=\beta I+K_{S,S}$. If a proposal $\widehat z$ has a verified residual bound
$\|b-A\widehat z\|_2\leq\eta$, then
\begin{equation}
 \|A^{-1}b-\widehat z\|_2\leq\frac{\eta}{\beta}.
 \label{eq:residual}
\end{equation}
This follows from $A\succeq\beta I$. For uncertain features the residual bound
must hold uniformly over their box, including directed product and summation
errors. Such bounds support coefficient enclosures and then the adaptive cell
checks; they do not make an unverified nominal solve exact.

Our prospective dispatcher admits the complete bounded solver schedule and its
explicit arrays before feature access. It prefers token space when $4T\leq d$
and streamed primal otherwise, subject to admission. This is a declared
engineering rule, not a learned optimal crossover. Compatible controls receive
the same kernels and choices. A trusted native exact Gram uses the same
feature-space certificate; its integer accumulator and archive parsing have
separate bounds.

\section{Conditional correctness and state}
The following statements concern the specified target and access contract. They
are not claims of universal completion, speed, cryptographic authenticity, or
novelty of interval verification.

\begin{theorem}[Conditional code correctness]
\label{thm:codes}
Assume that retained features are defined by the frozen source-local evaluator,
trusted descriptors satisfy~\eqref{eq:containment}, and every accepted rounding
certificate is valid for all feature realizations in its supplied box, conditional
on the accepted prefix. Assume that any fallback reproduces those same finite
features and certifies~\eqref{eq:recurrence}--\eqref{eq:rounding}. If the service
successfully commits all required stages, its model is $\Qfix(R)$.
\end{theorem}
\begin{proof}
Source locality makes each retained source's feature matrix equal to its
retained-only reconstruction. These matrices lie inside the prepared boxes.
For the first coordinate, a universally valid cell certificate therefore selects
the exact target code. Suppose the preceding codes equal the target prefix. The
certificate's conditional premise is then satisfied, so its next accepted code
also equals the target. Induction establishes every row. A successful fallback
uses the same features and exact point target, hence has the same conclusion.
The feature anchor is independent of calibrated outputs, so this argument applies
to every stage without changing a later stage's defining inputs. The complete
extent check establishes equality of the entire committed model.
\end{proof}

\paragraph{Canonical service state.}
Fix the codec, preparer, provenance policy, source ordering, target serialization,
and representation rule. Let $E(R)$ contain either retained canonical source
descriptors or the canonical exact pooled Grams required by that rule. Define
\begin{equation}
 \calS(R)=\operatorname{encode}\bigl(
   \text{target},\Qfix(R),R,E(R),\text{provenance}\bigr).
 \label{eq:state}
\end{equation}
Required tokens or source-retrieval obligations and external trust sidecars are
part of the service's storage/access accounting. The common pretrained checkpoint
and runtime dependencies remain required even if reported separately.
State provenance means retained-record provenance and target/codec bindings.
Execution logs and execution-provenance artifacts remain external; canonical
state equality alone does not compare them or the external Gram-trust sidecars.
The packed model artifact contains code indices for all 24 calibrated affine
stages; unchanged checkpoint components are still needed for inference.

\begin{lemma}[Exact source subtraction]
For fixed source-local features, $K(R)=\sum_{j\in R}X_jX_j^\top$. Exact addition and
subtraction of trusted source contributions, with original normalization
unchanged, produce the same rational Gram as fresh retained-source accumulation.
\end{lemma}
\begin{proof}
The summands are unchanged by other sources' deletion. Finite sums over disjoint
source sets are additive in exact rational arithmetic. Subtracting precisely the
removed summands therefore leaves exactly the retained sum. Canonical integer
scaling, source ordering, and serialization remove representation dependence on
addition or deletion order.
\end{proof}

\begin{theorem}[Canonical successor equality]
\label{thm:state}
Under Theorem~\ref{thm:codes}, suppose source descriptors are deterministic and
retained unchanged, exact Gram updates satisfy the preceding lemma, and all
compared executions use the same policies in~\eqref{eq:state}. A successful repair
returns the abstract canonical model and state $\calS(R)$, and therefore the same
bytes as any successful certified fresh retained-input reconstruction in the
same representation. This does not guarantee that a bounded fresh numerical
reconstruction completes. The equality conclusion holds after any
sequence of successful valid deletions.
\end{theorem}
\begin{proof}
The model agrees by Theorem~\ref{thm:codes}. Each retained descriptor agrees because
its source-local inputs and deterministic encoder agree. Each pooled Gram agrees
by exact subtraction and canonicalization. Membership, target, provenance, and
serialization policies are identical, so every serialized component agrees.
Applying this argument to each valid successor proves the sequence claim.
\end{proof}

State equality is representation-specific: a Gram state need not equal a
compressed-descriptor state. Changing the preparer or provenance format can also
change state bytes while leaving model codes unchanged. Implementation tests must
therefore compare independently reconstructed actual payloads, not only output
hashes or parser roundtrips. The theorem assumes precisely the premises those
checks are intended to exercise.

\section{Complete cost and access accounting}
A repair may reduce live storage while increasing request latency. Let $P_a$ be
initial preparation for representation $a$, and $C_{a,k}$ its complete cost on the
$k$th actual successor transition. For a declared initial-model convention,
\begin{equation}
 L_a(m)=P_a+\sum_{k=1}^{m}C_{a,k}
 \label{eq:lifetime}
\end{equation}
is its preparation-inclusive lifetime cost. A break-even observation requires
comparing these complete sums under matched access and output contracts. Repeating
alternative deletions from one unchanged root is not a successor-state sequence.
No theorem above bounds~\eqref{eq:lifetime} by a competitor's cost.

Request clocks must account for loading and validating the predecessor, source
access, exact accumulation, every certificate attempt, fallback, serialization,
complete output, and disposal. Shared checkpoint load, native compilation,
controller work, and independent audits need explicit placement rather than an
implicit zero. Nested timings cannot be added to their enclosing timer as if they
were disjoint. Refused runs retain their full observed cost and are not successful
speed observations.

An exact pooled Gram can be compact at large token counts, but deleting a record
requires that record's contribution. If individual contributions are not stored,
the service must retain access to deleted-source tokens or another way to
regenerate them. The pooled Gram and a membership digest alone do not supply that
information. Compressed fallback similarly requires retained-source access.

Storage includes model codes, descriptors or Grams, targets, tokens or retrieval
indexes, provenance, and required trust metadata. We report the shared checkpoint
separately and in a complete deployment total. An exported model already embedded
in a state must not be counted twice. Audit archives are separate physical data;
logical deletion from live calibration state does not erase those archives.

\section{Historical development evidence}
\label{sec:history}
The V39/V40 experiments use the pinned DistilGPT2 checkpoint and thirteen real
WikiText-2 training articles, each contributing 128 tokens. Original normalization
is $N_0=1664$, ridge is $1/100$, and the base-only row grids have four bits. The
requests retain 128, 768, or 1,536 tokens as specified by their campaign. They are
alternative requests from the same root, not independent roots or successive
deletions. Each arm/size has one development observation.

These experiments cover only the first complete QKV stage: 2,304 output rows,
width 768, and 1,769,472 codes. All seventeen V39 and twelve V40 outputs matched
the exact stage codes within their recorded comparisons. V40 additionally
compared every native source Gram against its Python reference and matched the
retained deletion/fresh Gram archives. Table~\ref{tab:history} presents selected
observations, retaining the important adverse outcomes; complete arm tables are
in the immutable project reports~\cite{v39,v40}.

\begin{table*}[t]
\centering
\small
\begin{tabular}{llrrrr}
\toprule
Campaign & Arm & Retained tokens & Seconds & Payload bytes & Replay fallback\\
\midrule
V39 & Cached token &128&1.493786&1,626,898&No\\
V39 & Cached primal &128&14.760016&1,626,898&No\\
V39 & Compressed 48-bit &128&2.096019&1,554,233&No\\
V39 & Compressed 32-bit &768&56.414186&3,455,531&Yes\\
V39 & Cached token &1,536&86.450110&9,205,544&No\\
V39 & Cached primal &1,536&31.323544&9,205,544&No\\
V39 & Gram deletion (Python) &1,536&22.674841&6,256,845&No\\
\midrule
V40 & Cached primal &768&23.020064&5,071,731&No\\
V40 & Streamed compressed 40-bit &768&25.387995&4,045,355&No\\
V40 & Native Gram deletion &768&30.136403&6,255,273&No\\
V40 & Cached primal &1,536&34.310553&9,205,544&No\\
V40 & Streamed compressed 40-bit &1,536&36.965634&7,152,668&No\\
V40 & Native Gram deletion &1,536&18.835432&6,256,852&No\\
\bottomrule
\end{tabular}
\caption{Selected historical single-stage development observations. Every row
matches the declared exact stage target; the 32-bit row includes its failed
certificate and fallback. Payloads include descriptors or a Gram, stage metadata,
and packed codes, but exclude the checkpoint and raw tokens. Arm times include
loading, replay, accumulation, certification, diagnostics, and output; shared
construction/builds and comparison/receipt costs are separately recorded. These
are neither complete-service nor lifetime measurements. Cross-campaign timing
ratios are not controlled treatment estimates.}
\label{tab:history}
\end{table*}

\paragraph{What the stronger control changes.}
At 768 retained tokens, V40's forty-bit compressed arm saves 20.2372\% of its
payload against cached primal, while taking 10.2864\% longer. At 1,536 tokens,
native Gram deletion takes 18.835432 seconds versus compression's 36.965634
seconds, and compression has 14.3174\% more payload. Omitting that stronger Gram
control would misrepresent the result. The useful possibility is a
workload-dependent storage--latency tradeoff, not universal compression speedup.

The V39 thirty-two-bit certificate refused at row 40, coordinate 424. Regenerating
the retained features and completing the registered fallback cost 31.454036
seconds within the 56.414186-second arm. This outcome illustrates why certificate
acceptance and fallback must be measured, rather than inferred from a nominal
compression ratio.

V40's native source accumulation took 9.397484 seconds versus 54.821483 seconds
for the Python reference, a component ratio of 5.833634. Preparation also ran both
implementations and compared them, so that ratio is not a complete preparation,
repair, or lifetime speedup. Its streamed primal implementation avoided retained
feature concatenation but was slightly slower in both observations; the explicit
working-set bound is not a measured whole-process memory saving.

\paragraph{Audit limits.}
The temporary runners compared actual binaries before shutdown. The fresh local
metadata audits checked published source identities, receipts, ledgers, and
ratios, but 77 derived V39 and 46 derived V40 binaries were absent locally. Those
metadata audits cannot independently repeat missing numerical comparisons.
Neither pilot measures language quality or supplies population confidence
intervals. Sixty previously exposed validation articles remain excluded from
future confirmation.

\section{Complete-model development pilot: V43}
\label{sec:v43}
The registered integration workload took the first three article IDs in the
published V39 development sampling order and the first 32 existing tokens from
each. The original target used $N_0=96$. It deleted the first sorted source ID,
retained 64 tokens, then deleted the next ID from that actual successor and retained
32 tokens. These short prefixes are exposed development inputs. They are chosen
to test complete-model integration and state transitions, not to establish useful
calibration scale or language quality.

The service covered all 24 stages and 42,467,328 codes. Separate arms used a lossless
cache, forty-bit compressed descriptors with exact fallback, and a declared hybrid
native representation. The hybrid stored native exact pooled Grams at the eighteen
width-768 stages and lossless factors at the six width-3072 MLP-down stages. It is
not a pure-Gram full-model experiment. Cold reconstruction built a complete
lossless state using the same optimized point dispatcher. A separate oracle
reexecuted retained source tokens and freshly constructed all three representations
without reading predecessor descriptors or Gram payloads. This oracle is
independent of predecessor evidence; it shares the point solver and codec
implementation, so it is not an independent implementation of the mathematics.

\paragraph{Actual code and state checks.}
All eight arm--successor comparisons passed the cluster byte audit: each complete
model equalled the oracle model, and each canonical state equalled freshly
constructed state in its representation. The copied artifacts also passed local
rehashing and direct model/state byte comparison. The first transition changed
2,774,109 code indices; the second changed 2,337,980. These counts compare adjacent
packed models under identical row grids, so the result is not merely retention of
an unchanged model. All 48 compressed stage certificates accepted without neural
replay or exact-feature fallback. Cached repair also replayed no sources; cold
reconstruction traversed 48 then 24 stage--source pairs. Hybrid deletion replayed
the deleted source through all 24 stages on each request before exact subtraction.

\begin{table*}[t]
\centering
\small
\begin{tabular}{llrrrr}
\toprule
Arm & Retained tokens & Request seconds & Live-state bytes & Deployment bytes & Peak MiB\\
\midrule
Lossless cache &64&36.161&37,599,262&390,424,437&940.6\\
Lossless cache &32&25.546&30,386,821&383,211,996&940.6\\
Compressed 40-bit &64&64.713&34,065,453&386,890,628&1,136.5\\
Compressed 40-bit &32&40.959&28,614,690&381,439,865&1,136.5\\
Hybrid native Gram &64&403.980&125,546,927&478,372,102&1,273.5\\
Hybrid native Gram &32&397.422&119,987,469&472,812,644&1,273.5\\
Cold lossless &64&86.999&37,599,262&390,424,437&918.2\\
Cold lossless &32&49.928&30,386,821&383,211,996&918.2\\
\bottomrule
\end{tabular}
\caption{Complete-model V43 development observations, one per actual successor
request. Live state includes packed model codes, both target manifests, retained
tokens/provenance, representation payloads, and the Gram-trust sidecar. Declared
deployment adds the shared 352,825,175-byte checkpoint/config. It excludes runtime
installation, audit archives, and separate execution-provenance receipts. Request
clocks cover predecessor loading/validation through complete output and disposal;
shared worker startup/admission is outside them. Peak RSS is measured for the
whole two-request worker, repeated in its two rows.}
\label{tab:v43}
\end{table*}

\paragraph{Cost result.}
Compression was 78.96\% and 60.34\% slower than lossless cached repair, while
reducing live state by 9.40\% and 5.83\%. With the common checkpoint counted, the
storage reductions were only 0.905\% and 0.462\%. Thus this pilot does not establish
a practically competitive compression advantage. The cache is faster but retains
more bytes; neither that distinction nor the particularly expensive hybrid at
these very small token counts supports universal dominance claims. Compression
was faster than cold reconstruction per request, which alone would be an
incomplete comparison.

Preparation used 22.783 seconds for context loading, 77.195 seconds for common
source extraction, and 39.994 seconds for the common original model. Representation
encoding, serialization, output, and disposal then cost 6.779 seconds for lossless,
44.130 for compression, and 46.618 for the hybrid. The combined preparation
transaction, which constructed all three representations, took 243.354 seconds.
Attributing the common preparation plus only each arm's own representation cost,
then adding its complete two-request controller transaction, gives 237.209 seconds
for cached repair, 318.162 for compression, 1,016.411 for hybrid repair, and 313.410
for cold reconstruction. Under this declared initial-model convention, compression
was 34.13\% more costly than the cache and 1.52\% more costly than cold
reconstruction. These are attributed totals, not separately measured deployments:
5.850 seconds of combined preparation overhead remains separately reported,
and no future break-even is extrapolated. The request-worker transactions include
their own context startup; the per-request columns in Table~\ref{tab:v43} do not.

\paragraph{Registration, resources, and audit boundary.}
The run used frozen source commit \texttt{304a9cc}, program \texttt{e350109a}, and
main audit \texttt{54cc9d3d}; complete identities and artifact bindings are retained
in the repository evidence~\cite{v43}. The registered runtime was Linux x86\_64,
CPython 3.12.14, with one institutional Slurm CPU, a 16-GiB virtual-address limit
and matching Slurm memory allocation. No GPU arithmetic defined the target.
All six corrected-campaign workers completed, with settled charges of 242, 90,
134, 826, 166, and 262 CPU seconds for preparation, cache, compression, hybrid,
cold, and oracle respectively. Together with the preserved 24-second failed
initial attempt, the total was 1,744 CPU seconds within the 13,000-second
continuation ceiling. That initial attempt failed in progress logging before
feature extraction; it was not a numerical refusal. Main audit analysis took
155.848 seconds separately from this worker ledger. Software and supplemental
audit costs are also separate.

The local archive retains 1,269,053,685 bytes of binary artifacts; its local
check rehashes all 57 declared binary and metadata artifacts and compares actual
model/state bytes. Local validation did not repeat model inference. A completed
supplemental audit (Slurm 13210, receipt \texttt{cd03e4c3}) additionally binds 138
input artifacts to the frozen program and sealed worker receipts. All eight
comparisons pass actual model/state byte equality, exact retained membership,
model-export/code consistency, cross-representation target/model agreement, and
registered-record binding. This audit took 102.184 wall seconds and 101.711 CPU
seconds, separately from the worker ledger. It checks retained artifacts and
provenance; it is not another numerical reconstruction or protection against a
hostile host. These checks support the narrow integration result above, not
scalable performance or a competitive compression system.

\paragraph{Exposed quality diagnostic (V44).}
A separately registered diagnostic evaluates the actual final V43 model against
matched sequential calibration, nearest rounding, and full precision~\cite{v44}.
Both calibrated targets use the final 32 retained tokens, original normalization
96, ridge $1/100$, and the same frozen row grids and checkpoint. The diagnostic
uses the same eight previously exposed V30 validation articles, truncated to 128
tokens each, with 1,016 next-token predictions per model. It consumes no new
quality inputs and preserves all sixty historical validation exclusions.

Before accepting new losses, all sixteen full-precision and nearest-rounding
article losses passed the frozen historical parity tolerance of $10^{-8}$ in
absolute mean negative log likelihood (NLL); the largest deviation was
$5.371\times10^{-15}$. Likelihoods use the shared NumPy binary64 evaluator. They
are approximate likelihood measurements, not certified inference or evidence of
GPU equivalence.

\begin{table}[t]
\centering
\small
\setlength{\tabcolsep}{3pt}
\begin{tabular}{lrrr}
\toprule
Model & Mean NLL & PPL & PPL/FP \\
\midrule
Full precision & 3.819909 & 45.60008 & 1.00000 \\
Nearest rounding & 4.172389 & 64.87023 & 1.42259 \\
Matched sequential & 3.966798 & 52.81515 & 1.15822 \\
Fixed anchor & 3.983710 & 53.71594 & 1.17798 \\
\bottomrule
\end{tabular}
\caption{V44 descriptive quality on eight already exposed articles. Perplexity
(PPL) is the exponential of pooled mean NLL; FP denotes full precision. Raw
per-article losses and all comparisons are retained in the completion receipt.}
\label{tab:v44}
\end{table}

The fixed/sequential pooled PPL ratio is 1.0170556 and the largest article ratio
is 1.0532114. Both frozen development guards (at most 1.05 pooled and 1.20 for
every article) pass. Fixed-anchor loss is nevertheless higher on six articles
and lower on two; its pooled PPL is 1.71\% worse than matched sequential and
17.80\% worse than full precision. The guard pass is a bounded development
diagnostic, not superiority, a population estimate, or untouched confirmation.
No inferential interval is reported. The one-use worker (Slurm 13215) completed
with 94 settled CPU seconds under its separate 1,202-second phase allowance.
Its source commit is \texttt{db95069}, program \texttt{e3d7b6df}, and completion
receipt \texttt{98870a81}; no adverse-result retry or threshold adjustment was
used.

\section{Related work and contribution boundary}
GPTQ develops post-training quantization using approximate second-order
information~\cite{gptq}. Our reconstruction contract must be stated independently:
finite source-local fixed-anchor features, original normalization, exact dyadic
interpretation, frozen grids, and lower midpoint ties. Agreement with this target
does not certify agreement with an unspecified GPTQ implementation or with a
sequential calibration target whose downstream activations change.

Exact additive sufficient statistics, ridge systems, low-rank identities,
residual-based verification, and margin stability are established methods. The
lemmas and conditional arguments above organize those ingredients around a
specific archived-evidence contract; their presentation is not a priority claim.
Any publishable contribution must identify what discrete recovery from genuinely
uncertain compressed features adds, and demonstrate useful costs against equally
optimized exact-cache and Gram controls.

\paragraph{Quantization-enabled exact deletion.}
Ginart et al.~\cite{ginart} retain metadata for quantized $k$-means, check whether
deletion changes quantized centroids, update metadata when stable, and otherwise
retrain. Their randomized-lattice analysis concerns expected deletion costs.
Thus stability checks, exact deletion aided by quantization, retraining fallback,
and unchanged models with updated metadata are established prior work.

Exact-Fun~\cite{exactfun} changes federated training to a quantized procedure,
retains historical models, corrects deleted gradients, checks recomputed
quantized parameters, and retrains from an affected round. Its full author-version
Theorem 2 proof uses uniform within-cell positions and independent coordinates;
the overlap calculation also uses a uniform perturbation model. Its probability
formula is not a distribution-free consequence of a displacement bound alone.
Our intended distinction is universal recovery of the specified adaptive codes
over uncertain feature evidence, including changed codes, rather than the
general fact that rounding can leave an output unchanged after deletion.

\paragraph{Executable certificates.}
ExecCert~\cite{execcert} separates candidate generation, the release contract, and
the verifier. It studies executable residual certificates, transfer to exact
references, guarded updates and rebuilds, decision agreement, and lifecycle costs.
Its Section 4.4 explicitly discusses finite output formats and codebooks. Our
adaptive weight-code contract over uncertain calibration features specializes a
different obligation from ridge-relative or fixed-readout tolerances; this does
not establish that the broader framework cannot express that contract.

These comparisons narrow the candidate contribution to exact adaptive-code
recovery from a lossy evidence representation, together with independently
checked canonical successor state. The conditional theorems alone do not
establish a new verification principle. A strong paper still needs a clearly
identified certificate or evidence advantage and useful complete costs. The
primary-source review is not an exhaustive priority search, and no execution or
reproduction of these external implementations is claimed.

\section{Remaining validation and limitations}
The next evaluation must separate three questions: does the service return the
specified model and state; does that model retain useful prediction quality; and
do its complete costs improve a meaningful deployment tradeoff? Exactness alone
answers only the first.

\paragraph{Scale and transfer.}
After the integration gate, use prospectively registered calibration sizes below,
near, and above stage width, independent corpus roots, and a second complete model.
The six width-3072 stages require their own admission and measured process costs.
A first-QKV timing cannot be multiplied by 24 to estimate full-model performance.
Measure actual changing-state sequences with matched representation choices and
all required source access.

\paragraph{Quality.}
Use untouched article-level evaluation units with overlap checks against exposed
data. Preserve all sixty historical validation exclusions; re-downloading the
same validation file does not make it unexposed. Compare with full precision,
nearest rounding, and an explicitly matched sequential-calibration reference.
Register useful endpoints and uncertainty analysis before evaluation. Do not tune
codec precision, solver routing, or stopping rules on future confirmation data.

\paragraph{Confirmation and mechanism.}
Freeze method choices before prospective randomized repetitions. Include
preparation, compilation, cache conditions, failures, replay, and serialization.
Evaluate compression precision, interval verification, solver dimension, and
source access through matched ablations. A single observation per arm and a
shared calibration root do not support a population confidence interval.

\paragraph{Deployment limits.}
The trusted computing base includes the finite evaluator, directed arithmetic,
compiler/runtime assumptions, source provenance, and immutable inputs during
verification. Hashes are integrity commitments under that trust model, not proof
against a hostile host. The current exact evaluator runs on supported CPU
arithmetic; using a GPU feature kernel would require a separate equivalence
argument or a separately declared target. Positive ridge does not guarantee
bounded numerical certification, and refusal must remain a visible outcome.

Logical successor-state deletion leaves audit artifacts and original external
corpora intact. It does not erase information in pretrained parameters, enforce
legal data-removal obligations, or certify physical memory or filesystem erasure.
Those are separate contracts. Most importantly, this draft remains a research
program with conditional correctness and bounded development evidence, not a
submission-ready claim of general efficient language-model unlearning.

\section{Reproducibility and evidence status}
Historical V39 and V40 reports identify their registration, source, workflow and
final evidence commits~\cite{v39,v40}. The V41 claim matrix fixes the interpretation
used here~\cite{boundaries}. Future releases should preserve every failed attempt,
source snapshot, resource ledger, selected-source commitment and actual comparison
receipt; derived binaries needed for independent reproduction must be distributed
or reproducibly regenerated under a new registration.

The completed V43 main, local, and supplemental artifact checks and the V44
exposed-quality diagnostic have separate receipts and resource accounts. The
listed external papers' titles, authors, and years were checked against primary
sources; this focused check does not replace a broader priority review.

\pending{Before sharing as a paper, add useful-scale, untouched-quality,
transfer, and prospective confirmation results, and complete the broader
comparative evaluation. Compile and visually inspect this source after the final
evidence edits.}

\begin{thebibliography}{9}
\bibitem{gptq}
Elias Frantar, Saleh Ashkboos, Torsten Hoefler, and Dan Alistarh.
2023. \emph{GPTQ: Accurate Post-Training Quantization for Generative Pre-trained
Transformers}. ICLR. Primary preprint:
\href{https://arxiv.org/abs/2210.17323v2}{arXiv:2210.17323, version 2}.

\bibitem{ginart}
Antonio A. Ginart, Melody Y. Guan, Gregory Valiant, and James Zou.
2019. \emph{Making AI Forget You: Data Deletion in Machine Learning}.
Advances in Neural Information Processing Systems 32.
\href{https://papers.nips.cc/paper_files/paper/2019/file/cb79f8fa58b91d3af6c9c991f63962d3-Paper.pdf}{Primary proceedings paper}.

\bibitem{exactfun}
Zuobin Xiong, Wei Li, Yingshu Li, and Zhipeng Cai.
2023. \emph{Exact-Fun: An Exact and Efficient Federated Unlearning Approach}.
IEEE International Conference on Data Mining, pages 1439--1444.
\href{https://doi.org/10.1109/ICDM58522.2023.00188}{Conference record};
\href{https://github.com/zuobinxiong/zuobinxiong.github.io/blob/300a0672be8d265ddebf20d6e69f790456a1bd81/assets/pdf/ExactFedUnlearning.pdf}{pinned full author version},
Theorem 2, pages 6--7. The proof comparison uses this historical author version.

\bibitem{execcert}
Ziyu Zhao, Xinyu Wang, Xiao-Wen Chang, and Yixuan He.
2026. \emph{From Mathematical to Executable Certificates for Machine Unlearning}.
\href{https://arxiv.org/html/2610.02268v1}{arXiv:2610.02268, version 1}.

\bibitem{v39}
Project \emph{second}. 2026. \emph{Real-data scale results and research decision:
V39}. Immutable project report at commit
\texttt{fbcc98c} (full identity in link):
\href{https://github.com/priyankjairaj100/second/blob/fbcc98c6559de4b88ffc9c8b960cf680d89e29d3/docs/SCALE_RESULTS_V39.md}{report and evidence pointers}.
Original completed workflow:
\href{https://github.com/priyankjairaj100/second/actions/runs/38023401448}{run 38023401448}.

\bibitem{v40}
Project \emph{second}. 2026. \emph{Native Gram and streamed scale results: V40}.
Immutable project report at commit
\texttt{fbcc98c} (full identity in link):
\href{https://github.com/priyankjairaj100/second/blob/fbcc98c6559de4b88ffc9c8b960cf680d89e29d3/docs/NATIVE_SCALE_RESULTS_V40.md}{report and evidence pointers}.
Original completed workflow:
\href{https://github.com/priyankjairaj100/second/actions/runs/38025361543}{run 38025361543}.

\bibitem{boundaries}
Project \emph{second}. 2026. \emph{Paper claim boundaries and evidence obligations}.
V41 handoff, immutable commit
\texttt{fbcc98c} (full identity in link):
\href{https://github.com/priyankjairaj100/second/blob/fbcc98c6559de4b88ffc9c8b960cf680d89e29d3/docs/CLAIM_BOUNDARIES_V41.md}{target, claims, and remaining obligations}.

\bibitem{v43}
Project \emph{second}. 2026. \emph{Complete fixed-anchor service development
pilot: V43}. Frozen source \texttt{304a9cc}:
\href{https://github.com/priyankjairaj100/second/blob/304a9cc73b3c263032f3ef6a02487b4d2df40fbe/research_v43/PROTOCOL.txt}{prospective protocol}.
Full program and main byte audit identities are in
\path{local_runs/full-model-v43-20261010-b/program.json} and
\path{local_runs/full-model-v43-20261010-b/artifact-audit.json};
the supplemental receipt is
\path{local_runs/full-model-v43-20261010-b/postrun-audit.json}
(\texttt{cd03e4c3}), and the retained-payload check is
\path{validation/v43_local_archive.json}.

\bibitem{v44}
Project \emph{second}. 2026. \emph{Bounded exposed-quality diagnostic: V44}.
Registered source \texttt{db95069}:
\href{https://github.com/priyankjairaj100/second/blob/db95069b2c3484ac615052fb1c1a711c8c773c2d/research_v44/PROTOCOL.txt}{prospective protocol}.
Full program identity is retained in
\path{local_runs/exposed-quality-v44-20261010-b/program.json}; results and
per-article losses are in
\path{local_runs/exposed-quality-v44-20261010-b/quality/outputs/completion.json}
(\texttt{98870a81}).
\end{thebibliography}
\end{document}
